Durante el desarrollo del presente proyecto se han identificado diversas oportunidades de mejora y líneas de evolución del sistema que no han sido abordadas en esta primera versión, pero que resultan coherentes con los objetivos y el alcance definidos en la introducción.

En particular, el sistema ha sido concebido con un enfoque modular y extensible, lo que permite la incorporación progresiva de nuevas funcionalidades y la mejora de las ya existentes sin necesidad de rediseñar la arquitectura base. Las líneas de trabajo futuro que se presentan a continuación se organizan en función de los módulos funcionales definidos en el sistema, reflejando así su evolución natural.

\subsection*{M1. Gestión de acceso, alta y autorización de usuarios}
\begin{itemize}
    \item Implementación de mecanismos de recuperación de contraseña y verificación de cuentas de usuario.
    \item Incorporación de medidas adicionales de seguridad frente a accesos indebidos, como la limitación de intentos de autenticación.
    \item Gestión avanzada de sesiones, incluyendo el control de sesiones activas y su posible revocación.
    \item Desarrollo de herramientas específicas para la visualización y análisis de los registros de acceso y acciones administrativas.
\end{itemize}

\subsection*{M2. Gestión comercial y clientes}
\begin{itemize}
    \item Incorporación de vistas históricas y métricas agregadas sobre visitas, actividad comercial y cumplimiento de rutas.
    \item Mejora de la replanificación diaria ante incidencias, cambios de franja o cancelaciones de última hora.
    \item Mayor automatización en la preparación de jornadas, reutilizando patrones de visitas y configuraciones recurrentes.
\end{itemize}

\subsection*{M3. Catálogo, productos y coloración}
\begin{itemize}
    \item Refuerzo del módulo con importación masiva de productos y mantenimiento asistido del catálogo.
    \item Incorporación de estados editoriales o flujos de borrador/publicación para contenidos comerciales y técnicos.
    \item Mejora de la búsqueda semántica, filtrado avanzado y reutilización de recursos de apoyo entre líneas relacionadas.
    \item Ampliación de la experiencia de coloración con comparativas, agrupaciones comerciales adicionales y material visual complementario.
\end{itemize}

\subsection*{M4. Pedidos, entregas y cobros}
\begin{itemize}
    \item Creación de pedidos a partir del catálogo ya implantado en \texttt{M3}, permitiendo seleccionar productos y cantidades sin exponer precios unitarios al cliente.
    \item Gestión del historial de pedidos por cliente y consulta de su estado dentro del ciclo comercial.
    \item Modelado de entregas, incidencias logísticas y productos pendientes de servir.
    \item Registro de cobros totales o parciales y construcción del estado de cuenta asociado a cada cliente.
\end{itemize}

\subsection*{Evolución técnica transversal}
\begin{itemize}
    \item Incorporación progresiva de pruebas automatizadas de integración y de interfaz para reducir regresiones manuales.
    \item Reorganización interna del código por \textit{feature}, especialmente en el área comercial, para facilitar la entrada de \texttt{M4} y módulos posteriores.
    \item Refuerzo de observabilidad, trazas de errores y utilidades de administración para despliegues más controlados.
\end{itemize}
